%----------------------------------------------------------------------------
\chapter{A rendszer specifikációja} \label{chapter2}
%----------------------------------------------------------------------------

\section{Felhasználói követelmények}

A felhasználói követelmények részletezik a felhasználói elvárásokat a Zebrafish rendszerrel szemben. Ez a fejezet határozza meg azokat a követelményeket, amelyeknek a rendszernek meg kell felelnie a célközönség igényeinek kielégítése érdekében.

A Zebrafish platform négy fő felhasználói csoportra összpontosít, mindegyiknek különböző igényeivel és munkafolyamataival:

\begin{enumerate}
  \item \textbf{Rendszergazdák} - akik az összes összetett, több Compose-projektből álló egységeket kezelhetik
  \item \textbf{Projektmenedzserek} - akik a Zebrafish platformon saját alkalmazásokat kezelnek
  \item \textbf{Fejlesztők} - akik több Compose-projektből álló alkalmazásokat fejlesztenek
  \item \textbf{Szervízfelhasználók} - akik automatizált módon hozzáférhetnek a Zebrafish platformhoz CI/CD rendszerekből
\end{enumerate}

\subsection{Rendszergazdák}

A rendszergazdák felelősek a Zebrafish operációs rendszer telepítéséért, karbantartásáért és felügyeletéért. Elvárásaik a stabilitásra, hatékonyságra és egyszerű kezelhetőségre összpontosítanak.

\begin{requirementstable}
  \caption{Felhasználói követelmények rendszergazdák számára} \\
  
  \requirementtitlewordwrap{Rendszertelepítés,\\image-kezelés} &
  \requirementslist{
    \item A rendszergazdáknak képesnek kell lenniük az egész rendszer összes konténer image-ének egyetlen paranccsal történő letöltésére, frissítésére
    \item A Compose-projektek első telepítése során az image-ek letöltésének gyorsnak kell lennie, kihasználva a párhuzamos letöltési lehetőséget
    \item A rendszer nem pazarolhatja a lemezterületet; a megosztott image-rétegeket csak egyszer szabad tárolni
    \item Az image-ek automatikus frissítésének támogatottnak kell lennie a teljes rendszerre vonatkozóan
  } \\

  \requirementtitlewordwrap{Rendszerfrissítés és karbantartás} &
  \requirementslist{
    \item A teljes rendszer frissítésének gyorsnak és atomi műveletnek kell lennie
    \item A frissítési folyamatnak egyszerűnek és könnyen kezelhetőnek kell lennie
    \item A rendszernek stabilnak kell maradnia egy frissítést követően
    \item Frissítés során nem szabad elvesztődnie a konfigurációs fájloknak
    \item Frissíŧés után a rendszer különleges beavatkozás nélkül kell működjön
  } \\

  \requirementtitlewordwrap{Rendszerfelügyelet és monitorozás} &
  \requirementslist{
    \item A rendszergazdáknak egyetlen paranccsal átfogó képet kell kapniuk a rendszer állapotáról
    \item Az állapotellenőrzésnek tartalmaznia kell, hogy mely konténerek futnak és melyek állnak
    \item A platformnak terminál alapú felhasználói felületet (TUI) kell biztosítania a rendszer vizuális áttekintéséhez
  } \\

  \requirementtitlewordwrap{Projekt- és konténerkezelés} &
  \requirementslist{
    \item A rendszergazdáknak képesnek kell lenniük egy teljes Compose-projekt leállítására egyetlen paranccsal
    \item Lehetőséget kell biztosítani az összes futó konténer azonnali újraindítására
    \item A független projektek nem indulhatnak újra feleslegesen, amikor egy másik, nem kapcsolódó projekt függőségei frissülnek
    \item A projektek közötti függőségeket a rendszernek automatikusan fel kell ismernie és kezelnie
    \item A rendszergazdáknak képesnek kell lenniük az alkalmazások naplóinak megtekintésére és kezelésére
    \item A rendszergazdák az összes projektet láthatják és módosíthatják a fájlrendszerből
  } \\

  \requirementtitlewordwrap{Hálózatkezelés} &
  \requirementslist{
    \item A hálózatokat a rendszernek proaktívan kell létrehoznia
    \item A hálózatok konfigurációját meg kell tudja adni a Compose-fájlokban
    \item A rendszergazdának nem kell manuálisan beavatkoznia a hálózatok létrehozásába
  } \\

  \requirementtitlewordwrap{Adatbiztonság és helyreállítás} &
  \requirementslist{
    \item A szerver adatainak automatikusan szinkronizálására lehetőséget kell adni
    \item Az adatok integritását biztosítani kell
    \item Katasztrófa esetén gyors helyreállítási mechanizmusnak kell rendelkezésre állnia
  }
\end{requirementstable}

\subsection{Projektmenedzserek}

A projektmenedzserek a Zebrafish platformot saját alkalmazásaik kezelésére használják. Fő feladatuk a projektek létrehozása, konfigurálása és karbantartása. Fontos megemlíteni, hogy a projektmenedzserek nem látnak rá más felhasználók Compose-projektjeire, a rendszernek biztosítania kell a projektek közti izolációt.

\begin{requirementstable}
  \caption{Felhasználói követelmények projektmenedzserek számára} \\

  \requirementtitlewordwrap{Projektek karbantartása} &
  \requirementslist{
    \item Szükséges, hogy saját projektjeiket láthatni tudják a fájlrendszerből
    \item Képesnek kell lenniük új Compose-projektek létrehozására a fájlrendszerből
    \item Terminálhozzáférést kell biztosítani a projektmenedzserek számára, hogy kezelhessék projektjeiket
    \item Elvárják, hogy a projektek egymással tudjanak kommunikálni a megosztott hálózatokon keresztül
    \item A projektmenedzsereknek nem láthatják más felhasználók projektjeit, csak a sajátjukat
  } \\

  \requirementtitlewordwrap{Erőforrásokhoz való hozzáférés} &
  \requirementslist{
    \item A projektmenedzserek saját projektjeiket hozzá kell tudják más projektmenedzserek hálózataihoz kapcsolni
    \item A hálózati névfeloldásnak működnie kell a projekt határain belül
    \item A projektmenedzsereknek képesnek kell lenniük a saját hálózataik konfigurálására
    \item A volumok egymás között nem szabad megoszthatóak legyenek, csak a saját projektjeikhez férhetnek hozzá
  }
\end{requirementstable}

\subsection{Fejlesztők}

A fejlesztők a Zebrafish platformot alkalmazásaik fejlesztésére, tesztelésére és futtatására használják. Számukra a legfontosabb a megszokott munkafolyamatok támogatása és a konténerek közötti zökkenőmentes kommunikáció.

\begin{requirementstable}
  \caption{Felhasználói követelmények fejlesztők számára} \\

  \requirementtitlewordwrap{Compose-kompatibilitás} &
  \requirementslist{
    \item A Compose-projektek felépítésének és konfigurációjának menete nem szabad eltérjen a megszokott Compose munkafolyamatoktól
    \item A fejlesztőknek nem kell új szintaxist vagy struktúrát tanulniuk a meglévő Compose fájljaik használatához
    \item Az összes standard Compose direktívának támogatottnak kell lennie (pl. services, networks, volumes)
  } \\

  \requirementtitlewordwrap{Multi-Compose projekt támogatás} &
  \requirementslist{
    \item A különböző Compose-projektekben futó konténereknek képesnek kell lenniük egymással kommunikálni
    \item A projektek közötti hálózati kommunikáció felépítésének zökkenőmentesnek kell lennie
    \item A szolgáltatások közötti névfeloldásnak működnie kell projekt határokon keresztül is
    \item A megosztott erőforrások (volumes, networks) nem kell projekteken keresztül használhatóak legyenek
  } \\

  \requirementtitlewordwrap{CLI integráció} &
  \requirementslist{
    \item A Zebrafish CLI-val készített konténerek kezelhetőek kell legyenek a containerd és nerdctl parancsokkal
    \item A CLI parancsok megszokott Docker Compose parancsokhoz hasonlóan működjenek
    }
\end{requirementstable}

\subsection{Szervízfelhasználók}

A szervízfelhasználók automatizált rendszerekből, például CI/CD pipeline-okból férnek hozzá a Zebrafish rendszerhez. Számukra a legfontosabb, hogy a rendszer könnyen automatizálható legyen.

\begin{requirementstable}
  \caption{Felhasználói követelmények szervízfelhasználók számára} \\

  \requirementtitlewordwrap{Automatizált hozzáférés} &
  \requirementslist{
    \item A szervízfelhasználóknak képesnek kell lenniük a Zebrafish rendszerhez parancssori eszközökön keresztül hozzáférni
    \item A rendszerre csatlakozni kell lehessen a CI/CD rendszerekből, például GitHub Actions-on keresztül
    \item Az automatikus hozzáférés képes kell legyen terminál alapú parancsokat futtatni
    \item A különböző compose parancsok nem szabad manuális beavatkozást igényeljenek
  } \\

  \requirementtitlewordwrap{Biztonságos hozzáférés} &
  \requirementslist{
    \item A hozzáférés biztonságos kell legyen: nem szabad illetéktelen felhasználók számára elérhetővé tenni
    \item A szervízfelhasználók csak a saját projektjeikhez férnek hozzá
  }
\end{requirementstable}

\section{Rendszerkövetelmények}

\subsection{Funkcionális követelmények}

Az alkalmazásnak olyan felülettel kell rendelkeznie, amely lehetővé teszi a felhasználók számára a projektek, konténerek és image-ek kezelését. A Zebrafish választott felülete egy parancssori felület (CLI - command line interface), amely parancsok sorát biztosítja a Docker Compose projektek, image-ek és konténerek kezeléséhez. A CLI-nak lehetővé kell tennie a felhasználók számára olyan általános feladatok elvégzését, mint a projektek indítása, leállítása és naplóinak megtekintése.

\begin{requirementstable}
  \caption{Funkcionális követelmények - Rendszerkezelés} \\

  \requirementtitlewordwrap{Rendszer-\\kompatibilitás} &
  \requirementslist{
    \item A Zebrafish rendszer kompatibilis kell legyen a \texttt{x86\_64} és \texttt{ARM64} architektúrákkal
    \item A Zebrafish kell támogassa a BIOS és UEFI indítást is
    \item A rendszer futtatható kell legyen a legelterjedtebb virtualizációs platformokon (KVM és VMware)
  }

  \requirementtitlewordwrap{Parancssori interfész (CLI)} &
  \requirementslist{
    \item A CLI-nek SSH-interfészen keresztül kell elérhetőnek lennie.
    \item Letölthető kell legyen az összes definiált Compose-projekt összes image párhuzamosan egy CLI parancs lefuttatásával 
    \item Megjeleníthető kell a rendszer általános állapota, beleértve a futó és leállított konténereket
    \item Frissíthető kell legyen a rendszer magja egyetlen parancs segítségével
    \item A rendszer különböző komponensei (pl. Crony, containerd) menedzselhetőek kell legyenek a CLI-on keresztül
    \item Az SSH-interfésznek támogatnia kell a Zebrafish operációs rendszerben elérhető összes többi parancsot is.
  } \\
 
  \requirementtitlewordwrap{Compose projektkezelés} &
  \requirementslist{
    \item A Compose-fájloknak kompatibilisnek kell lenniük a Compose V2 formátumával
    \item A hagyományos Compose-projektekhez hasonlóan a multi-Compose projekteket is konténerek létrehozásával és karbantartásával kell megvalósítani
    \item A Zebrafish által létrehozott konténereket a szabványos containerd és nerdctl parancsokkal is lehet kezelni.
    \item A parancsoknak támogatniuk kell a \texttt{-f} flag-et alternatív compose fájlok megadásához
  } \\

  \requirementtitlewordwrap{Image és tárolókezelés} &
  \requirementslist{
    \item Az image-kezelésnek kompatibilisnek kell lennie az OCI formátummal
    \item A rendszernek képesnek kell lennie a containerd image-ek kezelésére
    \item A rendszernek támogatnia kell a párhuzamos image letöltést a Tokio aszinkron futtatókörnyezet használatával
  }
\end{requirementstable}

\begin{requirementstable}
  \caption{Funkcionális követelmények - Hálózat és biztonság} \\

  \requirementtitlewordwrap{Hálózatkezelés} &
  \requirementslist{
    \item A rendszernek automatikusan létre kell hoznia a Compose-fájlokban definiált hálózatokat
    \item A CNI pluginokat (bridge, host-local, firewall, portmap, tuning) támogatni kell
    \item A multi-project hálózati kommunikációt külön hálózati névterekkel kell kezelni
    \item A szolgáltatások közötti névfeloldást DNS-sel kell biztosítani
  } \\

  \requirementtitlewordwrap{Biztonság} &
  \requirementslist{
    \item A rendszer magjának megváltoztathatatlannak kell lennie (immutable)
    \item SSH hozzáférés csak nyilvános kulcsos hitelesítéssel engedélyezett
    \item Port knocking mechanizmust kell implementálni az SSH védelmére
    \item DNS over HTTPS (DoH) kötelező az összes DNS lekérdezéshez
    \item HTTP/3 alapértelmezett protokoll a külső kommunikációhoz
  } \\

  \requirementtitlewordwrap{Adatkezelés} &
  \requirementslist{
    \item A ZFS fájlrendszernek biztosítania kell az adatintegritást
    \item Automatikus pillanatfelvételek készítése konfigurálandó időközönként
    \item Inkrementális replikáció támogatása távoli szerverekre
    \item A konfigurációs változtatások overlay fájlrendszeren keresztüli kezelése
  }
\end{requirementstable}

\subsection{Nem funkcionális követelmények}

Külső követelményekkel olyan feltételeket határozunk meg, amelyeknek a Zebrafish rendszernek meg kell felelnie, hogy illeszkedjen a szélesebb technológiai, jogi és társadalmi környezetbe.

\begin{requirementstable}
  \caption{Külső követelmények (external requirements)} \\

  \requirementtitlewordwrap{Szabványok és kompatibilitás} &
  \requirementslist{
    \item A rendszernek teljes mértékben kompatibilisnek kell lennie az OCI (Open Container Initiative) szabványaival
    \item A Compose fájloknak követniük kell a Docker Compose V2 specifikációt
    \item A container image formátumoknak támogatniuk kell az iparági szabványokat (Docker, OCI)
  } \\

  \requirementtitlewordwrap{Jogi megfelelőség} &
  \requirementslist{
    \item Minden szállított szoftverkomponens licensze kompatibilis kell legyen a rendszer nyílt forráskódú licenszével
    \item A rendszernek követnie kell az országok közti szoftverátadási törvényeket
  } \\

  \requirementtitlewordwrap{Interoperabilitás} &
  \requirementslist{
    \item A rendszernek képesnek kell lennie szabványos CI/CD rendszerekkel való integrációra
    \item Támogatnia kell a szabványos monitorozó rendszerekkel való integrációt (Prometheus, Grafana)
    \item Támogatnia kell a konténerfuttató platformokkal való együttműködést (containerd)
  }
\end{requirementstable}

\begin{requirementstable}
  \caption{Termékkövetelmények (product requirements)} \\

  \requirementtitlewordwrap{Hatékonyság} &
  \requirementslist{
    \item A lemezhasználat nem haladhatja meg a 256 MB-ot az alaprendszer esetében
    \item Az alap memóriahasználat nem haladhatja meg a 512MB-ot alapkonfigurációban
    \item A lokális hálózati névfeloldásnak 0.2 ms-nál rövidebb válaszidővel kell működnie
    \item A hálózati teljesítmény nem degradálódhat jelentősen a projektek számának növekedésével
  } \\

  \requirementtitlewordwrap{Gyorsaság} &
  \requirementslist{
    \item Az image-ek letöltése a lehető leggyorsabb kell legyen, kihasználva teljesen a hálózatot
    \item Megfelelően gyors hálózati kapcsolat esetén (1 Gbit) a rendszernek képesnek kell lennie 20 konténerkép frissítésére kevesebb mint egy perc alatt.
  } \\

  \requirementtitlewordwrap{Megbízhatóság} &
  \requirementslist{
    \item A rendszernek támogatnia kell legalább 30 egyidejűleg futó konténert
    \item A rendszernek képesnek kell lennie a konténerek automatikus újraindítására hiba esetén
    \item A rendszertől le kell lehessen kérni a futó konténerek állapotát és naplóit
    \item A frissítések után a rendszernek konzisztens és stabil állapotban kell maradnia
  }
\end{requirementstable}

\begin{requirementstable}
  \caption{Szervezési követelmények (organizational requirements)} \\

  \requirementtitlewordwrap{Integrációs követelmények} & 
  \requirementslist{
    \item A Zebrafish automatikusan integrálódik a meglévő ZFS storage rendszerekkel
    \item A Zebrafish-nek a meglévő eszközökkel, a Prometheusszal és a Grafanával nyomon követhetőnek kell maradnia
    \item A Zebrafish-nek képesnek kell lennie a meglévő CI/CD rendszerekkel való együttműködésre, mint például GitHub Actions
    \item A rendszer integrálható kell legyen SSH kapcsolatok révén (parancsok futtatásávl)
  } \\

  \requirementtitlewordwrap{Fejlesztési követelmények} &
  \requirementslist{
    \item A forráskód nyilvánosan elérhető GitHub repository-ban
    \item Folyamatos integráció (CI) minden commit után
    \item Dokumentáció szükséges a rendszer telepíŧéséről
  }
\end{requirementstable}